\documentclass[11pt,a4paper]{article}
\usepackage{rx-programming-style}

\begin{document}

\rxmaketitle{RX PHP Programming Style}
\rxdocumentmeta{PHP source code in RX/Raditex web applications and libraries}{1.0}{11 September 2026}

\tableofcontents
\clearpage

\section{1. Purpose}\label{1-purpose}

This document defines the PHP programming style used in RX projects.

The goals are readable code, predictable structure, explicit behaviour,
maintainability over cleverness, separation of responsibilities, safe
database and web programming, and consistency between projects.

The words \textbf{shall}, \textbf{should}, and \textbf{may} mean
required, strongly recommended, and optional respectively.

\section{2. Core principles}\label{2-core-principles}

\subsection{2.1 Clarity before brevity}\label{21-clarity-before-brevity}

Code shall be written for the programmer who must read and maintain it
later. Do not reduce line count merely to make code shorter.

Avoid compressed code:

\begin{lstlisting}[language=PHP]
if ($active) $value = $condition ? $first : $second;
\end{lstlisting}

Prefer explicit control flow:

\begin{lstlisting}[language=PHP]
if ($active)
{
    if ($condition)
    {
        $value = $first_value;
    }
    else
    {
        $value = $second_value;
    }
}
\end{lstlisting}

\subsection{2.2 One primary
responsibility}\label{22-one-primary-responsibility}

A class, function, or file should have one clear primary responsibility.

Typical RX separation:

\begin{lstlisting}
public page
    request handling and page composition

list class
    list query and list rendering

record class
    fetch/insert/update of one database record

form class
    form construction

generic widget
    reusable presentation logic
\end{lstlisting}

\subsection{2.3 Prefer composition}\label{23-prefer-composition}

In OO PHP, composition should normally be preferred to deep inheritance.
Inheritance is appropriate only when the relationship is natural and
remains easy to understand.

\subsection{2.4 Avoid speculative
abstraction}\label{24-avoid-speculative-abstraction}

Do not create a generic builder, metadata engine, factory hierarchy, or
base class merely to eliminate a few repeated lines. Small explicit
duplication is preferable to an abstraction that hides business
behaviour.

\section{3. Files and filenames}\label{3-files-and-filenames}

Filesystem names shall use lowercase words separated by hyphens.

Correct:

\begin{lstlisting}
rx-page.php
page-footer.php
page-footer-panel.php
inventory-item-record.php
inventory-item-edit-form.php
\end{lstlisting}

Avoid:

\begin{lstlisting}
rx_page.php
PageFooter.php
inventoryItemRecord.php
\end{lstlisting}

A PHP file should normally define one major reusable class.

Use paths relative to the current file:

\begin{lstlisting}[language=PHP]
require_once(__DIR__ . '/rx-widget.php');
\end{lstlisting}

Do not depend on the process working directory.

\section{4. PHP identifiers}\label{4-php-identifiers}

PHP identifiers shall use lowercase words separated by underscores.

Correct:

\begin{lstlisting}[language=PHP]
class page_footer
class inventory_item_record

function rx_html_escape(...)
function fetch_database_record(...)

$database_connection
$device_uuid
$submitted_values
$number_of_rows
\end{lstlisting}

Avoid camelCase:

\begin{lstlisting}[language=PHP]
class PageFooter
$databaseConnection
$deviceUuid
$numberOfRows
\end{lstlisting}

Names shall describe purpose. Avoid \passthrough{\lstinline!$d!},
\passthrough{\lstinline!$r!}, \passthrough{\lstinline!$q!},
\passthrough{\lstinline!$c!}, \passthrough{\lstinline!$x!},
\passthrough{\lstinline!$tmp!}, and similar abbreviations merely to save
typing.

Prefer:

\begin{lstlisting}[language=PHP]
$device_record
$database_result
$query_string
$database_connection
$temporary_filename
$object_definition
\end{lstlisting}

This rule also applies to loop variables:

\begin{lstlisting}[language=PHP]
foreach ($object_definitions as $object_definition)
{
    ...
}
\end{lstlisting}

Short established technical names such as
\passthrough{\lstinline!$sql!}, \passthrough{\lstinline!$html!},
\passthrough{\lstinline!$uuid!}, and \passthrough{\lstinline!$url!} are
acceptable when unambiguous.

Constants shall use uppercase underscore style:

\begin{lstlisting}[language=PHP]
const DEFAULT_PAGE_SIZE = 100;
const MAXIMUM_PAGE_SIZE = 500;
\end{lstlisting}

\section{5. Formatting and braces}\label{5-formatting-and-braces}

RX PHP uses an Ellemtel/Allman-like brace style.

\begin{lstlisting}[language=PHP]
if ($record === null)
{
    return null;
}
\end{lstlisting}

Functions:

\begin{lstlisting}[language=PHP]
function rx_html_escape(mixed $value): string
{
    return htmlspecialchars(
        (string)$value,
        ENT_QUOTES | ENT_SUBSTITUTE,
        'UTF-8'
    );
}
\end{lstlisting}

Classes:

\begin{lstlisting}[language=PHP]
class page_footer implements rx_widget
{
    private page_footer_panel $left_content;

    public function render(): string
    {
        ...
    }
}
\end{lstlisting}

Do not write:

\begin{lstlisting}[language=PHP]
if ($value === '') return;
\end{lstlisting}

Use:

\begin{lstlisting}[language=PHP]
if ($value === '')
{
    return;
}
\end{lstlisting}

Use four spaces per indentation level. Do not use tabs for source
indentation.

Break long calls and conditions over several lines. Use blank lines to
separate logical groups.

\section{6. Types}\label{6-types}

New standalone PHP modules should normally use:

\begin{lstlisting}[language=PHP]
<?php

declare(strict_types=1);
\end{lstlisting}

Existing projects may introduce strict types incrementally where
compatibility permits.

Functions and methods should use explicit parameter and return types
whenever practical:

\begin{lstlisting}[language=PHP]
public function fetch(
    string $device_uuid
): ?array
{
    ...
}
\end{lstlisting}

Nullability shall be explicit. Do not use empty string, zero, and
\passthrough{\lstinline!null!} interchangeably when they represent
different states.

\section{7. Object-oriented PHP}\label{7-object-oriented-php}

Create a class because it represents a coherent responsibility, not
merely because PHP supports classes.

Use the narrowest appropriate visibility. Do not make members public
merely for convenience.

Use interfaces when several implementations genuinely share a contract:

\begin{lstlisting}[language=PHP]
interface rx_widget
{
    public function render(): string;
}
\end{lstlisting}

Keep inheritance shallow. Prefer object composition and delegation.

Constructors should establish a valid object state. Avoid surprising
external side effects in constructors when a named method would make the
operation clearer.

\section{8. Functions and methods}\label{8-functions-and-methods}

Function names shall describe the operation.

Avoid:

\begin{lstlisting}[language=PHP]
h()
q()
do_it()
proc()
\end{lstlisting}

Prefer:

\begin{lstlisting}[language=PHP]
rx_html_escape()
build_query_string()
fetch_device_record()
validate_submitted_values()
\end{lstlisting}

A function should normally perform one logical operation.

Early returns are acceptable when they simplify validation and reduce
nesting.

Boolean functions should read naturally:

\begin{lstlisting}[language=PHP]
is_valid_uuid()
has_permission()
record_exists()
\end{lstlisting}

\section{9. Arrays}\label{9-arrays}

Use arrays where they naturally represent collections or simple
structured data.

Do not replace every explicit class with a large undocumented
associative array.

Prefer descriptive keys:

\begin{lstlisting}[language=PHP]
$submitted_values = [
    'name' => $name,
    'status' => $status,
    'device_uuid' => $device_uuid
];
\end{lstlisting}

\section{10. Database programming}\label{10-database-programming}

Database connection creation shall be centralized.

Data values shall be passed as parameters:

\begin{lstlisting}[language=PHP]
$database_result = pg_query_params(
    $database_connection,
    $select_sql,
    [$device_uuid]
);
\end{lstlisting}

Never concatenate request values into SQL.

When dynamic schema, table, or column names are genuinely required they
shall be validated and escaped with the proper PostgreSQL identifier
mechanism.

Long SQL shall be formatted for readability:

\begin{lstlisting}[language=PHP]
$select_sql = "
    SELECT
        device_uuid,
        name,
        status
    FROM
        rscada.device
    WHERE
        device_uuid = $1::uuid
";
\end{lstlisting}

Single-record database operations should normally live in record
classes. Public page files shall not contain large
\passthrough{\lstinline!SELECT!}, \passthrough{\lstinline!INSERT!}, or
\passthrough{\lstinline!UPDATE!} statements.

\section{11. Web input and output}\label{11-web-input-and-output}

Treat \passthrough{\lstinline!$\_GET!},
\passthrough{\lstinline!$\_POST!}, cookies, headers, and uploaded files
as untrusted input and validate them before use.

Dynamic text inserted into HTML shall be escaped using the equivalent
of:

\begin{lstlisting}[language=PHP]
htmlspecialchars(
    $value,
    ENT_QUOTES | ENT_SUBSTITUTE,
    'UTF-8'
);
\end{lstlisting}

Use a descriptive project helper such as
\passthrough{\lstinline!rx\_html\_escape()!} when appropriate.

Use \passthrough{\lstinline!http\_build\_query()!} for query strings.

Successful form updates should normally use Post/Redirect/Get.

Passwords, database passwords, API keys, and similar secrets shall not
be hard-coded in application source.

\section{12. Errors and exceptions}\label{12-errors-and-exceptions}

Errors shall not be silently ignored.

Exceptions are appropriate for invalid application configuration,
required service failures, or impossible internal states.

Expected user validation failures should normally be handled as
validation results rather than exceptions.

User-visible errors shall not expose credentials or unnecessary internal
details.

\section{13. Comments and
documentation}\label{13-comments-and-documentation}

Comments should explain why, assumptions, protocol requirements, or
non-obvious constraints.

Poor:

\begin{lstlisting}[language=PHP]
// Increment counter
$counter++;
\end{lstlisting}

Useful:

\begin{lstlisting}[language=PHP]
// Increment the generation when the hardware counter wraps so the
// complete sample identity remains monotonic.
$counter_generation++;
\end{lstlisting}

Public reusable classes and non-obvious interfaces should have concise
documentation.

\section{14. Security}\label{14-security}

At minimum:

\begin{itemize}
\tightlist
\item
  validate external input;
\item
  escape HTML output;
\item
  parameterize SQL values;
\item
  restrict file paths;
\item
  do not trust client-supplied roles or UUIDs;
\item
  perform authorization server-side;
\item
  keep secrets out of source code;
\item
  avoid production stack traces;
\item
  use secure session handling.
\end{itemize}

\section{15. Performance}\label{15-performance}

Do not optimize without evidence.

When performance matters:

\begin{enumerate}
\def\labelenumi{\arabic{enumi}.}
\tightlist
\item
  measure;
\item
  identify the actual bottleneck;
\item
  optimize that bottleneck;
\item
  measure again;
\item
  document any less-obvious implementation introduced for performance.
\end{enumerate}

\section{16. Example}\label{16-example}

\begin{lstlisting}[language=PHP]
<?php

declare(strict_types=1);

require_once(__DIR__ . '/rx-widget.php');

class text_widget implements rx_widget
{
    private string $text;

    public function __construct(string $text)
    {
        $this->text = $text;
    }

    public function render(): string
    {
        return htmlspecialchars(
            $this->text,
            ENT_QUOTES | ENT_SUBSTITUTE,
            'UTF-8'
        );
    }
}
\end{lstlisting}

\section{17. Review checklist}\label{17-review-checklist}

Before committing PHP code, check:

\begin{enumerate}
\def\labelenumi{\arabic{enumi}.}
\tightlist
\item
  Filenames are lowercase with hyphens.
\item
  PHP identifiers are lowercase with underscores.
\item
  Names are descriptive and single-letter abbreviations are avoided.
\item
  Braces and four-space indentation are consistent.
\item
  Each class/function has one coherent responsibility.
\item
  Composition is preferred over unnecessary inheritance.
\item
  External input is validated.
\item
  HTML output is escaped.
\item
  SQL data values are parameterized.
\item
  Database operations are in the proper layer.
\item
  Errors are handled explicitly.
\item
  Secrets are outside source code.
\item
  The code is clearer than the abstraction it replaces.
\end{enumerate}

\section{18. Core rule}\label{18-core-rule}

\begin{quote}
Make responsibilities, data flow, and control flow visible.
\end{quote}

Readable explicit code is preferred to clever compact code.


\end{document}
